%---------------------------------------------------------------------
\section{Conclusion}

Hangboarding is an evidence-supported supplemental training method for intermediate-to-advanced male rock climbers. The meta-analytic evidence (Stien et al.\ \cite{stien2023}) confirms consistent short-term improvements in dynamometric finger strength, RFD, and endurance (SDM 0.41--1.23), but no effect on climbing-specific on-wall performance. Max-hang protocols best develop MFS; repeater protocols best develop stamina and endurance; 80\% MFS protocols offer a pragmatic compromise. Daily low-intensity hangs (Abrahangs) did not significantly outperform climbing-only controls as a standalone protocol in the only available observational data ($p = 0.12$); they may add benefit when combined with max hangs but require RCT confirmation before being recommended broadly. Injury risk appears experience-dependent: experienced climbers may benefit from hangboarding's protective effect, while rapidly advancing less-experienced climbers face elevated injury risk. Critical gaps are the absence of female-cohort data, long-duration ($>\!10$~week) trials with injury tracking, and any trial demonstrating grade improvement. Until such data exist, practitioner folklore should be applied as biologically plausible guidance, not established protocol.
